\subsection{Finite Parameter Space Formalization}
\label{subsec:parameter_space}

To configure metaheuristics automatically using \texttt{irace}, we must convert the context-free grammar presented in Section~\ref{sec:design_space} into a finite parameter search space.
Because BNF grammars can generate infinite derivation trees via recursive calls or unbounded sequences, mapping them to \texttt{irace} requires bounding search depth and operator counts \citep{mascia2013grammars, mascia2014grammar, lopez2016irace}.

Following \citet{mascia2013grammars}, we bound recursion to a depth of two: a primary search procedure that can nest an inner search, where the inner search excludes \texttt{ils} and accepts only non-recursive algorithms (LS or TS).
Similarly, we limit multi-neighborhood search sequences in ILS to at most three distinct neighborhood structures, controlled by an integer parameter $\texttt{neighborhood\_number} \in \{1, 2, 3\}$.
For each active slot $j \in \{1, 2, 3\}$, \texttt{irace} configures the operator (\texttt{neighborhood}$_j$) and its exploration mode (\texttt{neighborhood\_exploration}$_j$), with \texttt{relax} conditioned to bypass exploration modes.
Categorical decisions, such as initial solution generation, priority dispatching rules, and sequence schedulers, map directly to categorical parameters.
Numerical and ordinal parameters governing Tabu Search, including \texttt{tenure}, \texttt{elite\_list\_size}, \texttt{limit\_iters}, \texttt{max\_d}, \texttt{max\_c}, and \texttt{decrease\_divisor}, are defined as conditional parameters that activate only when TS is selected.
Furthermore, a continuous parameter \texttt{max\_milli\_ils\_iteration} $\in [100, 30000]$ limits the execution time of each ILS iteration in milliseconds, preventing complex neighborhoods from exhausting the total time budget prematurely.
This time-bounding parameter allows \texttt{irace} to prevent computationally expensive neighborhood explorations from consuming an excessive portion of the total execution budget.

Because stochastic parameter generation can sample duplicate operators within a neighborhood sequence, an R repair hook intercepts candidates generated by \texttt{irace} before evaluation. The hook replaces repeated neighborhoods with distinct valid operators, ensuring neighborhood diversity across the sequence.

\subsection{Setup and Execution}
\label{subsec:irace_setup}

For the \texttt{irace} training phase, we generated 72 dedicated instances following the structural parameters and distributions of the standard benchmark \citep{baptiste2008lagrangian}, keeping the original test set disjoint.
To discard trivial problems, we solved each generated instance with CPLEX for 1 hour, retaining only those that remained unsolved to proven optimality.

We configured \texttt{irace} with a budget of 25,000 runs, setting a per-evaluation cutoff time of 60 seconds and parallelizing execution across 6 threads.
To ensure deterministic execution times, CPLEX calls within \texttt{relax} and \texttt{cplex} were restricted to a single thread (\texttt{threads=1}).

The framework was implemented in C++ and is openly available at \citet{leles2026jitjss}. Configuration runs used \texttt{irace} 4.4.3 \citep{lopez2016irace} in R, solving exact sub-problems via IBM ILOG CPLEX 22.1.1.0 \citep{cplex22}. Pseudorandom numbers were generated via the PCG library \citep{oneill:pcg2014} to guarantee statistical reproducibility across parallel runs.

Experiments were conducted on a dedicated machine featuring an Intel Core i5-11300H CPU @ 3.10 GHz, 24 GB RAM, and Arch Linux.
